\section{周：名分最后一次退出政治舞台}

\eventref{WQ-0403-THREEJIN}{战国·三家分晋}之后，洛阳的王廷还开着朝会，也还保有祭祀、册命和一套人人听得懂的礼仪语言；只是它已无力把这套语言变成命令。韩、赵、魏受册时，周王给出的仍是旧制度中最重的一层认可，三家实际掌握的却是城邑、赋役和军队。到齐、魏互称王，别国相继效仿，王号本身也从周王独有的名分变成列国竞逐的称号。残余周廷遂处在一个尴尬的位置：诸侯仍可借它的古老来装点自身的正当性，却不再愿意把争端、边界或出兵交给它裁断。

周室的衰微并不等于一夜之间只剩一块牌位。王畿内的权力早已分裂为西周、东周两支及其周边城邑；它们经营有限的土地、人口和与强国周旋的道路，王与两周君的关系也不总能从后出的叙事中清楚分辨。周赧王在位甚久，所能动用的却主要是名分、宗庙与洛邑这一地点本身。对于秦、韩等强邻而言，这一小块空间仍有价值：伊阙以北的通道、洛阳附近的城邑和一段可被用来联络诸侯的政治象征，都不能完全忽略。可是周廷已不能像春秋早期那样集合诸侯，也不能以王命调拨资源；它在列国外交中存活，更多靠各方暂时觉得保留它有用。

\subsection{一座王城，怎样分成两套小朝廷}

这份存活原本也有物质的一面。前五世纪中叶，周王室依洛阳周边田邑供养祭祀、朝会与宗族，\eventanchor{WQ-0438-ZHOU-ROYAL-ESTATES}{战国·周王室维持王畿田邑}。田邑的数量、收入和役夫没有可供核算的账簿，不能把王城想象成一座仍由完整财政支撑的都城；但宗庙、使者和王族若还要运转，便必有土地、粮食与人手被勉强留在这一小片空间。正因资源如此有限，王廷内部的分配不再只是宗室私事。

前367年前后，周室分出西周君与东周君两个政治中心，\eventanchor{WQ-0367-TWO-ZHOUS}{战国·周室分为东西周君}；传世叙事又把这件事系于周显王分置两周的安排，\eventanchor{WQ-0367-ZHOU-TWIN-COURTS}{战国·东西周封君中心形成}。分封的具体日期、边界和双方各得何种赋役都不清楚，不能画出精确界线；较可靠的含义是，洛阳附近的道路、使聘与收入从此可被两套封君网络分别经营。列国若要借道、游说、安置人质或联络旧臣，也不再只面对一个周廷。

周显王仍可用朝会与册命维持残余联系，\eventanchor{WQ-0346-ZHOU-XIAN-COURT}{战国·周显王维持王畿朝会}，但这更像把各方已经拥有的力量放进旧礼的字句，而非重新得到指挥它们的力量。前321年显王卒、慎靓王继位，\eventanchor{WQ-0321-ZHOU-XIAN-DEATH}{战国·周显王卒、慎靓王继位}，继承过程的记录极简，恰与周廷的处境相称：礼仪尚在，足以留下王名与继嗣；足以改变列国边界的军政资源却已在别处。

\begin{event}{公元前256年}{周赧王卒，东周王朝终结}{可靠纪年}{洛邑及其西侧城邑}\eventref{WQ-0256-ZHOUDIE}{战国·东周亡}
秦向东扩展时，西周君曾试图依附诸侯合纵，以伊阙方向为压力点牵制秦军。这是弱小政权在强邻逼近时几乎仅有的选择：若不联络东方诸国，便可能被逐一吞没；若联络失败，反会给秦以先发制人的口实。秦昭襄王五十一年，秦军攻取韩地阳城、负黍，西周君随即降服。传世纪年将秦取九鼎宝器、迁西周君于惮狐与周赧王卒连在这一年；留在洛邑的王廷至此失去最后可见的政治支点。铜鼎在这里不只是宝物：九鼎在古代叙事中承载王权的可见象征，秦把它们收走，正把一次对小国的处置写成了对旧共主秩序的接收。
\end{event}

这里的“东周亡”须作一层辨析。若“东周”指平王东迁以后延续的周王朝，前256年周赧王卒确是通常采用的终点；若指王畿分裂后由东周君控制的那个小政权，它在洛阳附近还延续到前249年\eventanchor{WQ-0249-EASTZHOU}{战国·东周君国亡}，才因与诸侯谋伐秦而被吕不韦处置、土地并入秦。故不能把前256年简单说成秦当年已消灭一切名为“东周”的实体，也不能因东周君尚存七年，便否认王朝意义上的周天子已不复存在。两种用法混在一起，正是晚周史最常见的年代误读之一。

\subsection{末局的账与鼎}

秦攻西周君而使其降服、献出城邑，\eventanchor{WQ-0256-QIN-XIZHOU}{战国·秦攻西周君}，是王朝终局的直接军事背景。秦接下的究竟是哪些田邑、仓廪和道路，史书没有一份移交清册；因此，\eventanchor{WQ-0256-QIN-ZHOU-HANDOVER}{战国·秦接管周王畿资源}只能说明这片政治空间开始接入秦的东进体系，不能把它写成当年已经完成的周人登记。西周君一旦屈服，王廷所剩的名分更难再靠邻国相助维持。

“债台高筑”的故事也应放回这一窘境。后世称赧王为合纵筹军费、向富户借贷，\eventanchor{WQ-0256-ZHOU-DEBT-MEMORY}{战国·周赧王债台高筑的记忆}与\eventanchor{WQ-0256-ZHOU-NANWANG-DEBT}{战国·周赧王为合纵筹费}，把财政枯竭凝成一个极有力量的寓言；债额、借贷程序与债主都无可靠账簿，不能把它当作可核的财政史料。它仍揭示了一个真实的限度：若一座王城想参与合纵，连筹集军费都须依赖城内富户，礼仪权威便已不能独自支撑对外行动。

九鼎的叙事因而格外醒目。前255年前后传世材料说周民东迁、九鼎入秦，\eventanchor{WQ-0255-NINE-TRIPODS-TO-QIN}{战国·周民东迁与九鼎入秦之说}；又把遗民、礼器与王畿田邑纳入秦的胜利记忆，\eventanchor{WQ-0255-QIN-ZHOU-RITUAL-CAPTURE}{战国·秦收取周室礼器的叙事}。九鼎是否果真按所述路线移走，没有考古实证，礼器去向亦不可一一坐实。较可以确定的是，周室人群确有分散，\eventanchor{WQ-0255-ZHOU-PEOPLE-EAST}{战国·西周亡后周民东迁}；旧官属、礼制与王统记忆遂成为东周君和列国仍可争取的象征资源，而非随王朝名义一同消失。

\subsection{最后的东周君}

东周君并未静待被并。前249年，他联络诸侯图抗秦，\eventanchor{WQ-0249-EASTZHOU-ANTI-QIN}{战国·东周君图合纵抗秦}，秦军遂攻而尽入其国。联军参加者和行动规模记载甚简，不能把最后一次尝试写成一场有完整盟书与阵图的大战；它却说明小政权仍试图用周统与道路位置换取外援。失败之后，秦处理周室的宗族、礼器与王畿遗存，\eventanchor{WQ-0249-DONGZHOU-ROYAL-RELICS}{战国·秦接收东周王室遗存}。物件清单和档案去向同样无从复原，但“收周”被塑造成继承王统的行动，正显示秦不仅夺取土地，也在预先占据统一天下所需的历史语言。

\begin{sourcenote}[王朝终点与残余政权]
\sourcebook{史记·周本纪}叙西周君惧秦而与诸侯约从，继而降秦、九鼎入秦、赧王卒；《秦本纪》可与其对读。《史记·秦始皇本纪》又记前249年东周君“与诸侯谋伐秦”，吕不韦“尽入其国”。《战国策·东周策》《西周策》保存两周君、游说与债务的后出记忆；《竹书纪年》系统为前255年周民东迁提供另一条编年线。这些都是传世编纂，不是前256年的现场记录；九鼎的迁移细节及两周君、王廷之间的具体权力关系，现代研究多结合洛阳地区考古与战国政治地理重新辨析。因此，王朝终结的年份较稳，行动者的每一步及其象征意义仍不宜写得过满。
\end{sourcenote}

真正消失的，是一个虽已虚弱、却仍可被所有列国援引的共同主位。它并没有直接造成\eventref{QIN-0230-UNIFICATION}{秦·统一六国}；秦的关中腹地、行政动员、战争经验，以及\eventref{WQ-0260-CHANGPING}{战国·长平之战}后愈发不对称的军力，才是统一得以推进的直接条件。但周王的退出改变了这些条件所处的政治语境：各国再不能把自己的主权安放在一个仍然存在的天子之下，合纵也少了一个可供共同调用的旧中心。秦后来面对的遂是一组以领土、军队和相王自立的国家；统一不再需要先解决“谁替周王主持天下”的问题，而转成“谁能把天下的征发、法令和道路纳入同一中心”的竞争。

\begin{turningpoint}[制度得失：名分留下的空位]
周廷残余曾给列国提供共同礼仪、交涉语言和最低限度的历史连续性，却没有足以承担仲裁、征发与防卫的资源。它的消失解除了一层旧秩序的象征约束，也使任何新中心都必须自行证明其统合能力。秦把这个空位填成皇帝、郡县和统一文书，并非周亡的自动后果；但一个真实共主不复存在，确实使那种答案更容易被想象，也更容易被强制推行。
\end{turningpoint}
